\begin{tikzpicture}[x=1cm,y=1cm,
  era/.style={rounded corners=4pt,fill=#1!7,draw=#1!60,line width=.7pt},
  eh/.style={anchor=north west,font=\sffamily\bfseries\fontsize{9.4}{11}\selectfont,text=#1!75!black,text width=3.6cm,align=left,execute at begin node={\hyphenpenalty=10000 }},
  et/.style={anchor=north west,font=\sffamily\fontsize{8.2}{9.6}\selectfont,text=Ink,text width=3.55cm,align=left,
    execute at begin node={\hyphenpenalty=10000 }}]
\def\W{3.82}
% ------------------------------------------------------------ era 1: search
\begin{scope}[shift={(0,0)}]
 \fill[era=DNAblue] (0,0) rectangle (\W,-7.7);
 \node[eh=DNAblue] at (.12,-.12) {1994--2003\\Search in a tube};
 \begin{scope}[shift={(.8,-3.15)},scale=.62]\Tube{(0,0)}{DNAsky}{}\end{scope}
 \fill[GelBg,rounded corners=1pt] (1.75,-1.65) rectangle (3.3,-3.2);
 \foreach \yy in {-1.95,-2.35,-2.8}{\fill[DNAblue] (2.15,\yy) rectangle (2.9,\yy+.07);}
 \node[et] at (.12,-3.45) {Libraries of candidate strands are filtered by PCR, gels and capture probes.\\[3pt]
   \textbf{Examples:} Hamiltonian path; satisfiability; maximal clique; sticker memory; a 20-variable 3-SAT instance.\\[3pt]
   \textbf{Lesson:} parallelism does not remove exponential cost (Part~I).};
\end{scope}
% ------------------------------------------------------------ era 2: machines and assembly
\begin{scope}[shift={(4.06,0)}]
 \fill[era=StrandRose] (0,0) rectangle (\W,-7.7);
 \node[eh=StrandRose] at (.12,-.12) {2001--2010\\Molecular machines};
 \Enzyme{(1.0,-2.35)}{.5}{\fontsize{6}{6}\selectfont cut}
 \draw[backbone,DNAblue,line width=1.6pt] (.35,-2.95)--(1.65,-2.95);
 \draw[backbone,StrandViolet,line width=1.6pt] (.35,-3.15)--(1.85,-3.15);
 \foreach \i in {0,...,3}{\foreach \j in {0,...,3}{
   \pgfmathtruncatemacro{\on}{mod(\i+\j,2)==0 && \i>=\j ? 1 : 0}
   \ifnum\on=1 \fill[StrandRose!60] (2.15+.36*\i,-1.65-.36*\j) rectangle +(.32,-.32);
   \else \draw[StrandRose!50] (2.15+.36*\i,-1.65-.36*\j) rectangle +(.32,-.32);\fi}}
 \node[et] at (.12,-3.45) {Enzymes and tiles act autonomously: automata step through states, tiles grow patterns, a scaffold folds into shapes.\\[3pt]
   \textbf{Examples:} enzyme-driven automata; Sierpinski crystals; DNA origami; enzyme-free logic.\\[3pt]
   \textbf{Lesson:} local rules compose (Parts~IV--V).};
\end{scope}
% ------------------------------------------------------------ era 3: dynamics
\begin{scope}[shift={(8.12,0)}]
 \fill[era=StrandGreen] (0,0) rectangle (\W,-7.7);
 \node[eh=StrandGreen] at (.12,-.12) {2006--2019\\Programmed dynamics};
 \draw[strand,StrandRose,line width=1.6pt] (.3,-1.75)--(1.6,-1.75);
 \draw[strand,StrandGreen,line width=1.6pt] (1.0,-2.15)--(2.4,-2.15);
 \draw[strandL,DNAblue,line width=1.6pt] (2.4,-2.5)--(.3,-2.5);
 \Rungs{1.1}{2.3}{-2.22}{-2.43}{.15}
 \draw[StrandAmber,line width=1.4pt,domain=0:1.1,samples=40,smooth] plot ({2.55+\x},{-3.1+.9/(1+exp(-9*(\x-.5)))});
 \node[et] at (.12,-3.45) {Strand displacement lets sequence program reaction rates: circuits, reaction networks, neural networks, walkers and robots.\\[3pt]
   \textbf{Examples:} seesaw circuits; Hopfield and winner-take-all networks; cargo-sorting robot; data storage.\\[3pt]
   \textbf{Lesson:} interfaces, leak and fuel (Parts~V--VI).};
\end{scope}
% ------------------------------------------------------------ era 4: learning and integration
\begin{scope}[shift={(12.18,0)}]
 \fill[era=EvidenceViolet] (0,0) rectangle (\W,-7.7);
 \node[eh=EvidenceViolet] at (.12,-.12) {2019--2026\\Learning, integration};
 \foreach \i/\y in {0/-1.6,1/-2.1,2/-2.6,3/-3.1}{\node[circle,fill=DNAblue!70,inner sep=0pt,minimum size=2.4mm] (in\i) at (.55,\y) {};}
 \foreach \o/\y in {0/-1.95,1/-2.75}{\node[circle,fill=StrandGreen!80,inner sep=0pt,minimum size=2.8mm] (out\o) at (1.7,\y) {};}
 \foreach \i in {0,...,3}{\foreach \o in {0,1}{\draw[EvidenceViolet!60,line width=.5pt] (in\i)--(out\o);}}
 \foreach \yy in {-1.6,-2.0,-2.4,-2.8}{\draw[Ink!60,line width=.5pt] (2.25,\yy)--(3.55,\yy);}
 \foreach \xx in {2.25,2.68,3.11,3.55}{\draw[Ink!60,line width=.5pt] (\xx,-1.6)--(\xx,-2.8);}
 \node[et] at (.12,-3.45) {Reprogrammable hardware, molecular learning, rechargeable circuits, recording in cells.\\[3pt]
   \textbf{Examples:} 21-algorithm tile set; gate arrays; learning in vitro; epi-bit storage; scaffolded computer.\\[3pt]
   \textbf{Lesson:} see the Epilogue.};
\end{scope}
% ------------------------------------------------------------ time axis
\draw[-{Stealth[length=2.6mm]},Ink,line width=1pt] (0,-8.05)--(16.3,-8.05);
\foreach \x/\y in {0.1/1994,4.16/2001,8.22/2006,12.28/2019}{\draw[Ink,line width=.8pt] (\x,-7.95)--(\x,-8.15);\node[slab,anchor=north] at (\x+.2,-8.15) {\y};}
\node[slab,anchor=north east] at (16.3,-8.15) {2026};
\end{tikzpicture}
